\section{Appendix - Spin Dependent Results}

WIMP-nucleon scattering can also have a spin-dependent interaction in which two limiting cases are considered: that WIMPs scatter only on protons or only on neutrons. Two isotopes of xenon have non-zero nuclear spin, \XeOneTwoNine\ (spin 1/2, \SI{26.4}{\percent} natural abundance) and \XeOneThreeOne\ (spin 3/2,  \SI{21.2}{\percent} natural abundance)~\cite{atomicweight}. As both have an unpaired neutron, the search is most sensitive to WIMP-neutron scattering. Sensitivity to a spin-dependent WIMP-proton interaction arises from mixing between proton and neutron spin states in isotopes with an unpaired neutron, albeit with increased uncertainty on the predicted signal rate~\cite{Engel:1991wq, klos2013structure, PhysRevLett.128.072502, PhysRevC.89.065501, PhysRevD.102.074018, Pirinen:2019gap}. 
Signal models for both the neutron and proton cases are constructed using the nuclear structure factors with uncertainties from Refs.~\cite{PhysRevD.102.074018,Pirinen:2019gap,PhysRevLett.128.072502}. This analysis quotes nominal limits which correspond to the mean structure functions from~\cite{PhysRevD.102.074018} and is chosen to facilitate a like-for-like comparison to previous limits from Xe-based experiments. An uncertainty is constructed for each $m_{\chi}$ by calculating the limit corresponding to the minimum and maximum interaction rate at each energy across the three models; this uncertainty also applies to the previous xenon results.
The details of data selection, background modeling, and statistical inference are identical to those reported in the main text.

\begin{figure}[htbp]
  \centering
  \includegraphics[ width=\linewidth]{figures/LZ_SR1_limit_SDn_PCL.pdf}
  \caption{The \SI{90}{\percent} confidence limit (black line) and uncertainty bands (gray) coming from xenon nuclear correction factors for the spin-dependent WIMP-neutron cross section vs.~WIMP mass using the mean of the nuclear structure factors from~\cite{PhysRevD.102.074018} and range  across~\cite{PhysRevD.102.074018,Pirinen:2019gap,PhysRevLett.128.072502}. Also shown are the PandaX-II~\cite{xia2019pandax}, LUX~\cite{akerib2017limits}, and XENON1T~\cite{collaboration2019constraining} limits. A similar uncertainty band as shown on this result applies to the other Xe-based limits.}
  \label{fig:limitplotSDn}
\end{figure}

\begin{figure}[htbp]
  \centering
  \includegraphics[ width=\linewidth]{figures/LZ_SR1_limit_SDp_PCL.pdf}
  \caption{The \SI{90}{\percent} confidence limit (black line) and uncertainty bands (gray) coming from xenon nuclear correction factors for the spin-dependent WIMP-proton cross section vs.~WIMP mass using the mean of the nuclear structure factors from~\cite{PhysRevD.102.074018} and range  across~\cite{PhysRevD.102.074018,Pirinen:2019gap,PhysRevLett.128.072502}. Also shown are the PICO-60~\cite{amole2019dark}, PandaX-II~\cite{xia2019pandax}, LUX~\cite{akerib2017limits}, and XENON1T~\cite{collaboration2019constraining} limits. A similar uncertainty band as shown on this result applies to the other Xe-based limits. The PICO-60 result relies on WIMP-scattering on the spin of the unpaired proton of $^{19}$F with minimal uncertainty.}
  \label{fig:limitplotSDp}
\end{figure}


Above the smallest tested WIMP mass of \SI{9}{Gev/c\squared}, the best-fit number of WIMP events is zero for both neutron and proton cases, and the data are thus consistent with the background-only hypothesis. Figure~\ref{fig:limitplotSDn} shows the 90\% confidence level nominal upper limit (black line) and nuclear structure function uncertainty on the limit (grey band) on the WIMP-neutron spin-dependent cross section as a function of mass. The minimum of the limit curve is at $m_{\chi} =$~\SI{30}{GeV/c\squared} at a cross section of $\sigma_{\mathrm{SD}}^{n} =$ \SI{1.49E-42}{cm\squared};  a power constraint is applied between \SI{13}{GeV/c\squared} and \SI{36}{GeV/c\squared}.  Figure~\ref{fig:limitplotSDp} shows the 90\% confidence level nominal upper limit and uncertainty on the WIMP-proton spin-dependent cross section as a function of mass. The minimum of the limit curve is at $m_{\chi} =$~\SI{32}{GeV/c\squared} at a cross section of $\sigma_{\mathrm{SD}}^p =$ \SI{4.2E-41}{cm\squared}; a power constraint is applied between \SI{13}{GeV/c\squared} and \SI{32}{GeV/c\squared}. The minimum and maximum limits which form the nuclear structure factor uncertainty are also power-constrained over the relevant mass range for both the neutron and proton cases.